%% ma: L10-B15-0047
%% lop: 10
%% chuong: 6
%% bai: 15
%% dang: 10.15.4
%% loai: TL
%% muc: TH
%% dap_an: "-"
%% nguon: "Ảnh giáo viên gửi 21/11/2026: Chuyên đề VI – Toán 10 – Chương VI (Hàm số – đồ thị và ứng dụng), tr. 199–203, Dạng 4 Câu 2"
%% kien_thuc: "Lớp 10 Bài 15 (hàm số đồng biến, nghịch biến)"
%% trang_thai: chinh-thuc
%% ghi_chu: "Dạng do giáo viên cung cấp (21/11/2026)."
\begin{ex}
Xét tính đồng biến và nghịch biến của hàm số $f(x)=\dfrac{x}{x-1}$ trên khoảng $(-\infty;1)$ và trên khoảng $(1;+\infty)$.
\loigiai{Với $x_1\ne x_2$ (khác $1$): $f(x_1)-f(x_2)=\dfrac{x_1}{x_1-1}-\dfrac{x_2}{x_2-1}=\dfrac{-(x_1-x_2)}{(x_1-1)(x_2-1)}$, nên $\dfrac{f(x_1)-f(x_2)}{x_1-x_2}=\dfrac{-1}{(x_1-1)(x_2-1)}$.

Trên mỗi khoảng $(-\infty;1)$ và $(1;+\infty)$, $x_1-1$ và $x_2-1$ cùng dấu nên tỉ số âm. Vậy hàm số nghịch biến trên từng khoảng.}
\end{ex}
